\chapter{Groupe abélien de type fini}
\lecture{7 octobre}

\section{Groupe abélien de type fini}
\begin{definition}[Groupe de type fini] \label{def: Groupe de type fini}
	On dit que les éléments $g_1, \ldots, g_k$ \textbf{engendrent} un groupe abélien $G$ si tout élément de $G$ peut être écrit comme combinaison linéaire 
	\[
	a_1 g_1 + \ldots + a_k g_k
	\]
	où $a_1, \ldots ,a_k \in \Z$. Si un groupe abélien admet un ensemble fini de générateur, alors on dira que $G$ est de \textbf{type fini}.
\end{definition}

\begin{proposition}
	Le groupe des nombres rationnels $\Q$ par rapport à l'addition n'est pas de type fini.
\end{proposition}

\prf{Supposons par l'absurde que $\Q$ sont de type fini. Alors il existe des rationnels
	\[
	\frac{b_1}{c_1}, \ldots , \frac{b_k}{c_k}
	\]
	tels que  pour tout éléments de $\Q$ puisse s'écrire comme
	\[
	a_1 \frac{b_1}{c_1} + \ldots + a_k \frac{b_k}{c_k}
	\]
	pour divers $a_1, \ldots, a_k \in \Z$. Mais le dénominateur doit diviser le nombre $c_1 \ldots c_k$, il est dont impossible que tout les nombres rationnels soient de cette forme.
}

\begin{definition}
	Pour tout groupe abélien $G$ et tout $n \in \Z$, la fonction
	\[
	G \to G, \quad g \mapsto ng
	\]
	est un morphisme. Le noyau de ce morphisme
	\[
	\Tors_n G = \ker(g \mapsto ng) = \big\{ g \in G \mid ng = 0 \big\}
	\]
	est appelé le $\bm n$-\textbf{ième sous-groupe de torsion}.
\end{definition}

\begin{proposition}
	Pour tout groupe abélien $G$, l'ensemble
	\[
	\Tors G = \bigcup_{i=1}^\infty \Tors_n G
	\]
	est un sous-groupe. Il sera appelé le \textbf{sous-groupe de torsion} de $G$. De plus un groupe abélien dont le sous-groupe de torsion est réduit au sous-groupe inverse est appelé \textbf{sans-torsion}.
\end{proposition}

\prf{Puisque $\Tors_n G$ est un sous-groupe de $G$ alors $\Tors G$ contient l'élément neutre et est fermé par inverse. De plus soit $g \in \Tors_n G$ et $h \in \Tors_m G$, alors
	\[
	nm(g+h) = nm g + nm h = m(ng) + n(mh) = 0
	\]
	ainsi $g+h \in \Tors_{mn} G\subset \Tors G$. Ainsi $\Tors G$ est bien un sous-groupe de $G$.
}

\begin{lemme}
	Pour tout groupe abélien $G$, le groupe quotient
	\[
	G / \Tors G
	\]
	est sans torsion.
\end{lemme}

\prf{Soit $g \in G$ tel que $n[g] = 0$ dans $G /\Tors G$ pour un certain $n \in \N$. Alors
	\[
	ng \in \Tors G
	\]
	et donc il existe $m \in \N$ tel que $mng = 0$. Ce qui implique que $g \in \Tors G$ et donc $[g] = 0$ dans $G/\Tors G$.
}

\begin{lemme}[Classification des groupes abéliens de type fini sans torsion] \label{lem: Classification des groupes abéliens de type fini sans torsion}
	Tout groupe abélien $G$ de type fini et sans torsion est isomorphe à un sous-groupe de $\Z^k$, pour un certain $k \geqslant 0$.
\end{lemme}

\prf{
Similairement à l'algèbre linéaire, on dira que $g_1, \ldots, g_k \in G$ sont \textbf{linéairement indépendants} si pour tout $a_1, \ldots, a_k \in \Z$,
\[
a_1 g_1 + \ldots + a_k g_k = 0 \implies a_1 = \ldots = a_k = 0
\]
Ainsi un considère un ensemble de générateur $g_1, \ldots, g_{k + l}$ de $G$, et supposons que $g_1, \ldots, g_k$ sont un sous-ensemble maximal d'éléments linéaire indépendants parmi les générateurs mentionnés. Ainsi, le morphisme
\[
\Z^k \overset\gamma\to G, \quad (a_1, \ldots, a_k) \mapsto a_1 g_1 + \ldots + a_k g_k
\]
est injectif. De plus par maximalité, on sait que pour tout $i \in \llbracket k + 1, l \rrbracket$, on peut écrire	
\[
b_i g_i = a_1^i g_1 + \ldots a_k^i g_k
\]
Soit $m$ le $\ppmc$ de $b_{k+1}, \ldots, b_l$, alors les $m g_i$ peuvent s'écrire comme combinaison linéaire des $g_1, \ldots, g_k$ pour tout $i \in \llbracket k + 1, l \rrbracket$. Et donc pour tout $g \in G$, on conclut que $mg$ peut s'écrire comme combinaison linéaire de $g_1, \ldots, g_k$. Et donc le morphisme
\[
G \overset\delta\to G,\quad g \mapsto mg
\]
est contenue dans l'image du morphisme injectif	$\gamma$. Cela implique que $\delta$ se factorise comme
\[
\delta : G \to \Z^k \overset\gamma\to G
\]
Cependant, $\delta$ est injectif car $G$ est sans torsion et donc le morphisme qu'on vient de construire $G \to \Z^k$ est aussi injectif. Ainsi, $G$ est isomorphe à l'image de ce morphisme.
}

\begin{lemme}[Classification des sous-groupes de $\Z^k$] \label{lem: Classification des sous-groupes de Zr}
	Tout sous-groupe de $\Z^k$ est isomorphe à $\Z^r$ pour un certain $r \geqslant 0$.
\end{lemme}

\prf{On résonne par récurrence, si $k = 1$, on sait que tout les sous-groupe de $\Z$ sont de la forme $n \Z$. Or il est claire que
	\[
	n \Z \simeq \Z
	\]
	
	Désormais, par récurrence, supposons que tout sous-groupe de $\Z^{k-1}$ soit abélien, soit $G \subset \Z^k$ un sous-groupe, alors on construit la suite exacte courte
	\[
	0 \to \Z^{k-1} \overset\gamma\to \Z^k \overset\pi\to \Z \to 0
	\]
	où le morphisme $\pi$ est la projection sur la dernière coordonnée. De plus on définit 
	\[
	K = \ker(\pi_{\mid G}) = \{(x_1, \ldots, x_k) \in G \mid x_k = 0\} \et L = \pi(G)
	\]
	ainsi cela induit naturellement la suite exacte courte
	\[
	0 \to K \overset{\gamma'}\to G \overset{\pi'}\to L \to 0
	\]
	Alors $L$ est un sous-groupe de $\Z$ puisque c'est l'image d'un isomorphisme, ainsi on a deux possibilités
	\begin{enumerate}
		\item Si $L = 0$, alors $G \simeq K$ est un sous-groupe de $\Z^{k-1}$, que l'hypothèse montre être abélien et libre.
		\item Si $L \simeq \Z$,  il nous reste à construire une section ou une rétraction, pour ce faire il suffit de prendre une section de la projection, pour ce faire on peut envoyer $1 \in L$ sur un élément quelconque dans $\pi'^{-1}(1)$. Ainsi, on a construit une suite exacte courte scindée de groupe abélien, et donc on a 
		\[
		G \simeq K \times L \simeq \Z^r \times\Z = \Z^{r+1}
		\]
	\end{enumerate}
}

\begin{corollaire} \label{cor: suite exacte courte scindée}
		Toute suite exacte courte de groupes abéliens
		\[
		0 \to K \to G \overset\pi\to \Z^l \to 0
		\]
		est scindée, et en particulier, $G \simeq K \times \Z^l$.
\end{corollaire}

\prf{Par l'algèbre linéaire, on a que le $e_i$ sont les générateur standard de $\Z^l$. Et puisque le morphisme $\pi : G \to \Z^l$ dans la suite exacte courte est surjectif, alors on peut choisir $g_i \in \pi^{-1}(e_i)$ pour tout $i$. Alors la fonction
	\[
	\psi : \Z^l	\to G, \quad (n_1, \ldots, n_l) \mapsto n_1 g_1 + \ldots + n_l g_l
	\]
	est un morphisme. De plus par construction cela défini bien un section au sens où $\pi \circ \psi = \mathrm{Id}_{\Z^l}$. Par conséquent la suite exacte courte est scindée.
}

\begin{theoreme} \label{thm: groupe abélien isomorphe à Zr fois Tosr G}
	Tout groupe abélien de type fini on a
	\[
	G \simeq  \Z^r \times \Tors G
	\]
	pour un certain $r \geqslant 0$.
\end{theoreme}

\prf{Par le \autoref{lem: Classification des groupes abéliens de type fini sans torsion}, on a une suite exacte courte
	\[
	0 \to \Tors G \overset\pi\to L \to 0
	\]
	où $L$ est un groupe abélien sans torsion. Mais comme $G$ est de type fini, alors $L$ l'est aussi. Et donc par le \autoref{lem: Classification des groupes abéliens de type fini sans torsion} et le \autoref{lem: Classification des sous-groupes de Zr}, on trouve que $L$ est isomorphe à $\Z^r$ pour un certain $r \geqslant 0$. Et finalement par le \autoref{cor: suite exacte courte scindée}, on peut conclure que
	\[
	G \simeq L \times \Tors G \simeq \Z^r \times \Tors G
	\]
}

\lecture{14 octobre}	

\section{Groupe de torsion $\Tors G$}

\begin{proposition}[Sous-groupe de type fini] \label{prop: Classification des sous-groupes de groupe de type fini}
	Si $G$ est un groupe de type fini, alors tout sous-groupe $H \leqslant G$ est également de type fini.
\end{proposition}


\prf{On procède par récurrence, supposons que la \autoref{prop: Classification des sous-groupes de groupe de type fini} soit vraie pour tout groupe abélien avec moins de $n$ générateurs. Soit $H \leqslant G$ un sous-groupe arbitraire. De plus soit $G' \leqslant G$ le sous-groupe engendrée par $g_1, \ldots, g_{n-1}$, et considérons la suite exacte courte
	\[
	0 \to G' \overset\imath\to G \overset\pi\to G/G' \to 0
	\]
	Similairement à la preuve du \autoref{lem: Classification des sous-groupes de Zr}, on pose
	\[
	K = \ker(\pi_{\mid H}) = H \cap G'  \et L = \pi(G)
	\]	
	ce qui induit la suite exacte courte
	\[
	0 \to K \overset{\imath'}\to H \overset{\pi'}\to L \to 0
	\]
	Comme $K$ et $L$ sont des sous-groupes, l'hypothèse de récurrence implique que $K$ et $L$ sont tout deux de type fini. Soit $x_1, \ldots, x_k$ des générateurs $K$ et $y_1, \ldots, y_l$ des générateurs de $L$. Alors
	\[
	\imath'(x_1), \ldots, \imath'(x_k), g_1, \ldots, g_l
	\]
	sont des générateurs de $H$, pour tout choix de $\{g_i \in \pi'^{-1}(y_i)\}_{i \in \llbracket1,l \rrbracket}$. En effet, pour tout $h \in H$, on peut écrire
	\[
	\pi'(h) = b_1 y_1 + \ldots + b_l y_l
	\]
	Ce qui implique que 
	\[
	h - b_1 g_1 - \ldots - b_l g_l \in \ker \pi' = \Im \imath'
	\]
	Mais comme $\imath'$ est injective alors par le \autoref{thm: Premier théorème d'isomorphisme} on a
	\[
	K \simeq \Ima \imath'
	\]
	et donc il doit exister $a_1, \ldots, a_k \in \Z$ tels que 
	\[
	h - b_1 g_1 - \ldots - b_l g_l = a_1 \imath'(x_1) + \ldots \imath'(x_k)
	\]
	Et donc la proposition suit.
}

\begin{definition} \label{def: produit et somme direct}
	Soit $G_1, G_2, \ldots$ une collection au plus dénombrable de groupes abéliens, alors on définit
	\[
	\prod_{i=1}^\infty G_i = \{(g_1, g_2, \ldots) \mid g_i \in G_i\} \et \bigoplus_{i=1}^\infty G_i =  \{ (g_1, g_2, \ldots) \mid g_i \in G_i, \et \exists k \in \N, \, \forall i \geqslant k, \quad g_i = 0\}
	\]
	Qu'on appelle respectivement le \textbf{produit direct} et la \textbf{somme directe}.
\end{definition}

\begin{definition}[Sous-groupe de $p$-torsion]
	Soit $G$ un groupe abélien quelconque et $p$ un nombre premier, alors on définit le \textbf{sous-groupe de $p$-torsion} comme
	\[
	A_p(G) = \bigcup_{n=1}^\infty \Tors_{p^n} G
	\]
	Il est facile de voir que $A_p(G)$ définit bien un sous-groupe.
\end{definition}

\begin{proposition} \label{prop: Decomposition du groupe de torsion}
	Pour tout groupe abélien, on a 
	\[
	\Tors G \simeq \bigoplus_\text{$p$ premier} A_p(G)
	\]
	avec la somme directes des groupes abéliens définie comme dans la \autoref{def: produit et somme direct}.
\end{proposition}

\prf{On définit naturellement le morphisme
	\[
	\varphi : \bigoplus_{i=1}^\infty A_p(G) \to \Tors G, \quad (g_{p_1}, g_{p_2}, \ldots) \mapsto \sum_{i=1}^\infty g_{p_i}
	\]
	Cette somme est bien définie puisque les $g_{p_i}$ sont nuls à partir d'un certain rang et donc on a une somme finie. Il reste donc à montrer que $\varphi$ définit en fait un isomorphisme.
	\begin{itemize}
		\item[$\bullet$] Injectif: Soit $g_{p_1}, \ldots, g_{p_k} \in G$ tels que
		\[
		p_1^{d_1} g_{p_1} = \ldots = p_k^{d_k} g_{p_k} = 0
		\]
		où $p_1, \ldots, p_k$ sont des nombres premiers distincts. Si la collection de ces $g_{p_i}$ est dans le noyau de $\ker \varphi$, c'est-à-dire
		\[
		g_{p_1} + \ldots + \ldots g_{p_k} = 0
		\]
		Alors en multipliant cette somme par $p_2^{d_2} \cdots p_k^{d_k}$, on trouve
		\[
		p_2^{d_2} \cdots p_k^{d_k} g_{p_1}+ \underbrace{p_2^{d_2} \cdots p_k^{d_k} g_{p_2}}_{=0} + \ldots + \underbrace{p_2^{d_2} \cdots p_k^{d_k} g_{p_k}}_{=0} = 0 \implies 	p_2^{d_2} \cdots p_2^{d_2} \cdots p_k^{d_k} g_{p_1} = 0
		\]
		Cependant, on sait aussi que $p_1^{d_1}g_{p_1} = 0$. Comme $p_1^{d_1}$ et $p_2^{d_2} \cdots p_k^{d_k}$ sont premiers entre eux, alors cela implique que $g_1 = 0$. Par un argument analogue, on montre que $g_1 = g_2 = \ldots = g_k = 0$, ce qui implique que le noyau $\ker \varphi$ est trivial est donc $\varphi$ est bien injectif.
		
		\item[$\bullet$] Surjectif: Soit $g \in \Tors_n G \subset \Tors G$ et considérons la décomposition en facteurs premiers
		\[
		n = p_1^{d_1} p_2^{d_2} \cdot p_k^{d_k}
		\]
		où $p_1, \ldots, p_k$ sont des nombres premiers distincts. Alors l'élément
		\[
		g_i = p_1^{d_1} \cdots p_{i-1}^{d_{i-1}} p_{i+1}^{d_{i+1}} \cdot p_k^{d_k}g
		\]
		a un ordre divisant $p_i^{d_i}$, et donc appartient à $A_{p_i}(G)$. Cependant les entiers naturels
		\[
		\Big\{  p_1^{d_1} \cdots p_{i-1}^{d_{i-1}} p_{i+1}^{d_{i+1}} \cdot p_k^{d_k}\}_{i\in \llbracket1,k \rrbracket}
		\]
		ont un $\pgcd$ égal à $1$. Et donc par Bézout, il existe des entiers $a_1, \ldots, a_k$ tels que
		\[
		\sum_{i=1}^k  a_i p_1^{d_1} \cdots p_{i-1}^{d_{i-1}} p_{i+1}^{d_{i+1}} \cdot p_k^{d_k} = 1
		\]
		Et donc par conséquent on a
		\[
		a_1 g_1 + \ldots + a_k g_k = \Bigg(\sum_{i=1}^k  a_i p_1^{d_1} \cdots p_{i-1}^{d_{i-1}} p_{i+1}^{d_{i+1}} \cdot p_k^{d_k} \Bigg) g = g
		\]
		ce qui implique que $g$ appartient à l'image du morphisme $\varphi$.
	\end{itemize}
}	

\begin{definition}[$p$-groupe]
	Soit $G$ un groupe, alors on dira que $G$ est un $p$-\textbf{groupe} si l'ordre de chaque élément de $G$ est une puissant de $p$.
\end{definition}

\begin{lemme} \label{lem: ordre d'un groupe abélien}
	Si un nombre premier $p$ divise l'ordre d'un groupe abélien $G$, alors $G$ est à un élément d'ordre $p$.
\end{lemme}

\prf{On va résonner par récurrence sur l'ordre de $G$. Premièrement si $G$ cyclique, alors le résultat est évident. Sinon, on peut considérer un élément $h \in G$ qui n'engendre par tout $G$. Si $p \mid a = \ord_G(h)$, il est facile de trouver un multiple de $h$ d'ordre $p$. Si $p \not \mid a$, on considère le sous-groupe $H$ engendré par $h$ et on considère le quotient
	\[
	G/H
	\]
	Comme l'ordre de $H$ est premier avec $p$, alors $p$ divise l'ordre de $G/H$. Par hypothèse de récurrence, il existe un élément $g \in G$ tel que $[g]$ est d'ordre $p$ dans $G/H$. Ce qui implique que
	\[
	pg \in H \implies apg = 0
	\]
	dans $G$. Et donc ce qui implique que $ag$ est soit d'ordre $p$, soit que $ag = 0$ Cependant, puisque $\pgcd(a,p) =1$, alors $pg \in H$ et $ag = 0 \in H$ impliquerait que $g \in H$, ce qui contredirait le fait que $[g]$ est d'ordre $p$ dans $G/H$.
}

\begin{theoreme} \label{thm: Classification des p-groupe}
	Un groupe abélien fini $G$ est un $p$-groupe si et seulement si $|G|$ est une puissance de $p$.
\end{theoreme}

\prf{Conséquence immédiate du théorème de Lagrange et du \autoref{lem: ordre d'un groupe abélien}.}


\begin{proposition} \label{prop: Classification des p-groupe par isomorphisme}
	Tout $p$-groupe abélien fini est isomorphe à
	\[
	\Z/p^{d_1}\Z\times \dots \times \Z/p^{d_k} \Z
	\]
	pour divers entiers positifs $d_1, \ldots, d_k$.
\end{proposition}

\prf{On résonner par récurrence sur $|G|$, on considère un élément de $G$ d'ordre maximal possible, on l'appelle $h$ et on suppose qu'il est d'ordre $p^{d_1}$. Soit $H \simeq \Z/p^{d_1} \Z$ le sous-groupe de $G$ engendré par $h$. Ainsi le groupe quotient
	\[
	G/H
	\]
	est un $p$-groupe en raison du \autoref{thm: Classification des p-groupe}. Et donc par hypothèse récurrence, il existe des entiers positifs $d_2, \ldots, d_k$ tels que
	\[
	G /H \simeq \Z/p^{d_2}\times \dots \times \Z/p^{d_k} \Z
	\]
	De plus comme tout élément de $G/H$ a un ordre inférieur dans $G$ et que $p^{d_1}$ soit l'ordre maximal, on conclut que $d_1 \geqslant d_2, \ldots, d_k$. Ce qui implique l'existence d'une suite exacte courte
	\[
	0 \to \Z/p^{d_1} \Z \to G \overset\pi\to \Z/p^{d_2}\times \dots \times \Z/p^{d_k} \Z \to 0
	\]
	Il nous reste donc à exhiber une section puisque tout les groupes en jeu sont abéliens. Pour ce faire, on définit pour tout $i \in \llbracket 2, k \rrbracket$,
	\[
	g_i \in \pi^{-1}(e_i)
	\]
	Alors on a
	\[
	p^{d_i}[g_i] = 0 \in G/H \implies p^{d_i} g_i = a_i h, \quad \forall i \in \llbracket 2, k \rrbracket
	\]
	pour un certain $a_i \in \Z$. On décompose $a_i = p^{s_i} t_i$ de sorte que $t_i$ soit premier avec $p$. Ce qui donne
	\[
	p^{d_i}g_i = p^{s_i}t_i h
	\]
	Mais comme tout les éléments de $G$ on a un ordre qui est par hypothèse inférieur à $p^{d_1}$, alors
	\[
	0 = p^{d_1} g_i = p^{d_1 - d_i + s_i} t_i h
	\]
	Puisque $t_i$ est premier avec $p$, alors l'ordre de $h$ est le même que l'ordre de $t_ih$ et donc nécessairement $d_1 - d_i + s_i \geqslant d_1 \implies s_i \geqslant d_i$ pour tout $i$. Donc on peut revenir à notre expression de toute à l'heure
	\[
	p^{d_i}g_i = p^{s_i}t_i h \implies p^{d_i}\underbrace{(g_i - p^{s_i-d_i}t_i h)}_{g_i'} = 0
	\]
	La formule ci-dessus garantit que l'application
	\[
	\Z/p^{d_2}\times \dots \times \Z/p^{d_k} \Z  \overset\psi\to G, \quad (x_2,\ldots, x_k) \mapsto x_2 g_2' + \ldots + x_k g_k'
	\]
	est un morphisme bien défini. Il est également clair que $\pi \circ \psi = \mathrm{Id}$ par construction.
}

\begin{theoreme}[Classification des groupes abélien de type fini]
	Tout groupe abélien de type fini $G$ est isomorphe à un produit direct
	\[
	\Z^r \times \Z/p_1^{d_1} \times \ldots \times \Z/p_k^{d_k} \Z
	\]
	pour un certain $r \geqslant 0$, appelé le \textbf{rang} de $G$ et divers puissances premières $p_1^{d_1}, \ldots, p_k^{d_k}$, appelés les \textbf{diviseurs élémentaires} de $G$. De plus la décomposition est unique à permutation des facteurs près.
\end{theoreme}

\prf{Par le \autoref{thm: groupe abélien isomorphe à Zr fois Tosr G}, tout groupe abélien de type fini est de la forme
	\[
	\Z^r \times \Tors G
	\]
	De plus par la \autoref{prop: Classification des sous-groupes de groupe de type fini} et puisque tout groupe abélien de torsion de type fini est fini, alors le sous-groupe de torsion en question peut s'écrire comme somme direct fini de $p$-groupes d'après la \autoref{prop: Decomposition du groupe de torsion}. Et finalement la \autoref{prop: Classification des p-groupe par isomorphisme} cela assure que tout ces facteurs sont des produits de la forme $\Z/p^d\Z$ pour divers $p$ premier.
}